\documentclass[11pt]{article}
\usepackage[margin=1in]{geometry}
\usepackage[T1]{fontenc}
\usepackage[utf8]{inputenc}
\usepackage{lmodern}
\usepackage{booktabs}
\usepackage{array}
\usepackage{longtable}
\usepackage{graphicx}
\usepackage{xcolor}
\usepackage{hyperref}
\usepackage{amsmath}
\hypersetup{
  colorlinks=true,
  linkcolor=blue,
  urlcolor=blue,
  citecolor=blue
}

\title{HMR3D Current Version Summary\\Implementation and Experimental Results}
\author{Repository snapshot: commit \texttt{46f5e73}}
\date{April 11, 2026}

\begin{document}
\maketitle

\section{Scope}
This document summarizes the currently committed HMR3D implementation in
\texttt{/home/pengyue/Codespace/HMR3D}. It intentionally covers only the
committed \texttt{pose-anchor} version and excludes later local
\texttt{Gaussian/submap} experiments that were not kept in the clean working
tree.

The current system extends the upstream \texttt{TTT3R} runtime with an explicit
state-level memory lifecycle:
\begin{itemize}
\item archive recurrent state snapshots into an external bank,
\item retrieve candidate memories by descriptor similarity,
\item verify candidates with a counterfactual geometric rollout, and
\item further filter ambiguous or stale recoveries with anchor-pose scoring.
\end{itemize}

The best validated configuration is the experiment manifest
\texttt{configs/tum\_relpose\_sweep\_224\_v11\_pose\_anchor\_full7.json}, evaluated
against the \texttt{TTT3R} baseline on seven TUM long-sequence windows under
deterministic inference.

\section{Implementation Summary}
\subsection{Code Organization}
The current implementation is split across the following local modules:
\begin{itemize}
\item \texttt{hmr3d\_memory/config.py}: memory hyper-parameters and mode-specific
configuration.
\item \texttt{hmr3d\_memory/router.py}: archive bank, retrieval logic, recovery
proposal construction, geometric verification statistics, and anchor-pose
verification statistics.
\item \texttt{hmr3d\_memory/adapter.py}: runtime wrapper around the upstream
\texttt{TTT3R} rollout, including the counterfactual verification path.
\item \texttt{hmr3d\_memory/eval\_relpose.py}: deterministic rel-pose
evaluation, per-sequence metric aggregation, and memory event dumps.
\item \texttt{scripts/run\_relpose\_memory\_sweep.py}: experiment sweep entrypoint.
\item \texttt{scripts/run\_longseq\_eval.py}: full long-sequence evaluation
entrypoint.
\end{itemize}

\subsection{Memory Lifecycle}
At the current commit, the memory system is state-level and does not modify the
upstream submodule parameters.

\paragraph{Archive.}
Every active segment maintains frame descriptors and the current recurrent state.
When the archive interval is reached, the system stores:
\begin{itemize}
\item a segment descriptor,
\item a short tail descriptor history,
\item a combined state descriptor from \texttt{state\_feat} and \texttt{mem},
\item the recurrent tensors needed to reconstruct the state, and
\item an optional camera pose anchor.
\end{itemize}
This is implemented in \texttt{MemoryRouter} via \texttt{ArchiveEntry}.

\paragraph{Retrieve.}
At each eligible frame, the router ranks archived segments with descriptor
similarity. It then applies thresholding on:
\begin{itemize}
\item query-to-archive descriptor similarity,
\item current-state similarity,
\item the query-state gap, and
\item a recovery cooldown and an attempt cooldown.
\end{itemize}

\paragraph{Geometric verification.}
For shortlisted candidates, the adapter runs a counterfactual rollout: it
injects the archived state proposal, executes the current frame, and compares
the resulting geometric quality against the baseline rollout. The candidate is
accepted only if it satisfies the configured geometric and confidence tests.

\paragraph{Anchor-pose verification.}
The current best version adds a second verifier for ambiguous or stale
candidates. It compares the current pose estimate and the archived camera anchor
through an anchor-pose quality score. The recovery is retained only if the
candidate score is at least as good as the baseline score under the configured
ratio and gain thresholds.

\subsection{Current Default Hyper-parameters}
The best validated manifest uses the following memory settings:
\begin{center}
\begin{tabular}{ll}
\toprule
Parameter & Value \\
\midrule
\texttt{archive\_interval} & 120 \\
\texttt{retrieval\_topk} & 2 \\
\texttt{retrieval\_similarity\_thresh} & 0.92 \\
\texttt{verification\_similarity\_thresh} & 0.80 \\
\texttt{max\_state\_similarity\_for\_recover} & 0.995 \\
\texttt{query\_state\_gap\_thresh} & -0.05 \\
\texttt{recovery\_alpha} & 0.35 \\
\texttt{verification\_geo\_ratio\_thresh} & 1.0 \\
\texttt{verification\_conf\_tolerance} & 0.1 \\
\texttt{min\_frames\_before\_retrieve} & 30 \\
\texttt{retrieval\_cooldown} & 30 \\
\texttt{retrieval\_attempt\_cooldown} & 30 \\
\texttt{reset\_on\_archive} & false \\
\texttt{enable\_anchor\_pose\_verification} & true \\
\texttt{anchor\_pose\_only\_for\_ambiguous} & true \\
\texttt{anchor\_pose\_rotation\_weight} & 0.25 \\
\texttt{anchor\_pose\_score\_ratio\_thresh} & 1.0 \\
\texttt{anchor\_pose\_min\_score\_gain} & 0.0 \\
\bottomrule
\end{tabular}
\end{center}

\section{Evaluation Protocol}
\subsection{Dataset and Hardware}
The current report uses the deterministic rel-pose sweep at:
\begin{itemize}
\item dataset: \texttt{tum\_s1\_1000},
\item input size: \texttt{224},
\item device: \texttt{cuda},
\item seed: \texttt{0},
\item deterministic: \texttt{true}.
\end{itemize}

The seven evaluated sequence windows are:
\begin{itemize}
\item \texttt{sitting\_static},
\item \texttt{static\_head500},
\item \texttt{static\_tail500},
\item \texttt{sitting\_xyz},
\item \texttt{xyz\_head500},
\item \texttt{xyz\_mid500},
\item \texttt{xyz\_tail500}.
\end{itemize}

The evaluation files are stored under
\texttt{reports/generated/tum\_relpose\_sweep\_224\_v11\_pose\_anchor\_full7}.

\subsection{Compared Methods}
Three configurations were compared:
\begin{itemize}
\item \textbf{TTT3R baseline}: upstream behavior without external archive/retrieve/recover.
\item \textbf{HMR geometry-only}: descriptor retrieval plus geometric verification.
\item \textbf{HMR pose-anchor}: geometry-only plus anchor-pose verification on
ambiguous or stale candidates.
\end{itemize}

\section{Main Results}
\subsection{Aggregate Results}
Table~\ref{tab:aggregate} reports the main aggregate metrics.

\begin{table}[h]
\centering
\caption{Aggregate results on the seven-window TUM rel-pose benchmark. Lower is
better for ATE, RPE\_trans, and RPE\_rot. Higher is better for FPS.}
\label{tab:aggregate}
\begin{tabular}{lrrrr}
\toprule
Method & Avg. ATE & Avg. RPE$_{\text{trans}}$ & Avg. RPE$_{\text{rot}}$ & Avg. FPS \\
\midrule
TTT3R baseline & 0.088360 & 0.025146 & 0.938209 & 3.779 \\
HMR geometry-only & 0.081347 & 0.025555 & 0.938447 & 3.716 \\
HMR pose-anchor & \textbf{0.079694} & 0.025560 & \textbf{0.934255} & 3.668 \\
\bottomrule
\end{tabular}
\end{table}

Relative to the baseline, the current \textbf{HMR pose-anchor} version:
\begin{itemize}
\item reduces average ATE from 0.088360 to 0.079694,
\item yields a relative ATE improvement of approximately 9.8\%,
\item slightly improves average rotational error, and
\item incurs an average throughput drop of about 2.9\%.
\end{itemize}

Relative to the geometry-only version, pose-anchor further improves ATE by
approximately 2.0\%.

\subsection{Per-sequence Results}
Table~\ref{tab:perseq} shows the per-window ATE comparison between the baseline
and the current best HMR version.

\begin{table}[p]
\centering
\caption{Per-sequence ATE comparison. Positive ``Rel. gain'' means the HMR
version improves over the baseline.}
\label{tab:perseq}
\begin{tabular}{lrrrr}
\toprule
Sequence & Baseline ATE & HMR ATE & $\Delta$ATE & Rel. gain (\%) \\
\midrule
\texttt{sitting\_static} & 0.028312 & 0.021613 & -0.006700 & 23.66 \\
\texttt{static\_head500} & 0.023047 & 0.021860 & -0.001187 & 5.15 \\
\texttt{static\_tail500} & 0.014938 & 0.013395 & -0.001544 & 10.33 \\
\texttt{sitting\_xyz} & 0.278476 & 0.269411 & -0.009065 & 3.26 \\
\texttt{xyz\_head500} & 0.064023 & 0.054727 & -0.009297 & 14.52 \\
\texttt{xyz\_mid500} & 0.067648 & 0.070329 & 0.002681 & -3.96 \\
\texttt{xyz\_tail500} & 0.142074 & 0.106521 & -0.035553 & 25.02 \\
\bottomrule
\end{tabular}
\end{table}

The key pattern is stable:
\begin{itemize}
\item 6 of 7 windows improve in ATE,
\item the largest gain appears on \texttt{xyz\_tail500},
\item the only regression is \texttt{xyz\_mid500}, and
\item the method is improving global trajectory quality more clearly than local
translational error.
\end{itemize}

\subsection{Memory Activity}
The best HMR trial generated the following aggregate memory statistics:
\begin{itemize}
\item 33 archives,
\item 108 retrieval attempts,
\item 32 successful recoveries after verification,
\item 124 geometric verification rollouts,
\item 32 anchor-pose verification accepts,
\item 92 anchor-pose verification rejects,
\item average recovery success rate of 29.6\%.
\end{itemize}

These numbers show that the gain does not come from blindly restoring history.
The majority of retrieval candidates are rejected by the verification layers,
and only a filtered subset is merged back into the active state.

\section{Interpretation}
The current version supports the following conclusions.

\paragraph{What works.}
The archive/retrieve/recover lifecycle can improve long-sequence robustness
without retraining the base model. The main benefit is on \textbf{global ATE},
which is consistent with the goal of mitigating drift and exploiting long-range
revisits.

\paragraph{What does not improve yet.}
Average translational RPE does not improve. This means the method is not yet a
uniform improvement across all metrics. It is currently stronger as a
\textbf{global drift correction mechanism} than as a local-frame refinement
mechanism.

\paragraph{Why pose-anchor helps.}
The geometry-only variant already removes many bad recoveries, but ambiguous
revisits still remain. Adding anchor-pose verification narrows those cases and
reduces average ATE from 0.081347 to 0.079694 on the same full-seven benchmark.

\section{Reproducibility}
The exact experiment manifest for the current best result is:
\begin{quote}
\texttt{configs/tum\_relpose\_sweep\_224\_v11\_pose\_anchor\_full7.json}
\end{quote}

The sweep can be reproduced with:
\begin{quote}
\texttt{python3 scripts/run\_relpose\_memory\_sweep.py --config}\\
\texttt{configs/tum\_relpose\_sweep\_224\_v11\_pose\_anchor\_full7.json}
\end{quote}

The key output files are:
\begin{itemize}
\item \texttt{reports/generated/tum\_relpose\_sweep\_224\_v11\_pose\_anchor\_full7/leaderboard.json}
\item \texttt{reports/generated/tum\_relpose\_sweep\_224\_v11\_pose\_anchor\_full7/ttt3r\_baseline/summary.json}
\item \texttt{reports/generated/tum\_relpose\_sweep\_224\_v11\_pose\_anchor\_full7/hmr\_full\_geometry\_pose\_anchor/summary.json}
\end{itemize}

\section{Current Limitations}
This committed version still has clear boundaries:
\begin{itemize}
\item it is a state-level memory system, not yet a map-level submap system,
\item it improves ATE more consistently than RPE$_{\text{trans}}$,
\item one of the seven windows still regresses, and
\item the evaluation was run at input size 224 because the local RTX 2060 6GB
machine could not reliably run the same long-sequence setup at 512.
\end{itemize}

\section{Bottom Line}
For the current clean repository state, the correct short summary is:

\begin{quote}
The committed HMR3D implementation adds an external state-memory lifecycle to
TTT3R and achieves a \textbf{9.8\% average ATE improvement} over the TTT3R
baseline on a deterministic seven-window TUM long-sequence benchmark, with a
small throughput cost and no corresponding gain in translational RPE.
\end{quote}

\end{document}
